% !TeX program = pdflatex
% !TeX encoding = utf8
\documentclass[12pt,oneside,openright,a4paper,brazil]{abntex2}

\usepackage{lmodern}
\usepackage[T1]{fontenc}
\usepackage[utf8]{inputenc}
\usepackage{indentfirst}
\usepackage{microtype}
\usepackage{graphicx}
\usepackage{float}
\usepackage{booktabs}
\usepackage{array}
\usepackage{multirow}
\usepackage{enumitem}
\usepackage{xcolor}
\usepackage{hyperref}
\usepackage{caption}
\usepackage{subcaption}
\usepackage{listings}
\usepackage{listingsutf8}
\usepackage{mathtools}

\definecolor{codegray}{gray}{0.95}
\definecolor{codeborder}{gray}{0.8}
\definecolor{primaryblue}{RGB}{0,62,116}

\setlength{\parindent}{1.25cm}
\setlength{\parskip}{0.15cm}

\newcommand{\includeplaceholder}[2]{%
  \IfFileExists{#1}{\includegraphics[width=\textwidth]{#1}}{%
    \fbox{\parbox[c][6cm][c]{0.9\textwidth}{\centering #2\\[0.5cm]\footnotesize Substituir por \texttt{\detokenize{#1}}}}}}

\newcommand{\code}[1]{\texttt{\detokenize{#1}}}

\lstdefinestyle{codestyle}{
  backgroundcolor=\color{codegray},
  frame=single,
  rulecolor=\color{codeborder},
  basicstyle=\ttfamily\small,
  keywordstyle=\color{primaryblue},
  commentstyle=\color{green!40!black},
  stringstyle=\color{orange!70!black},
  numbers=left,
  numberstyle=\tiny,
  stepnumber=1,
  numbersep=10pt,
  tabsize=2,
  showstringspaces=false,
  breaklines=true,
  inputencoding=utf8,
  columns=flexible
}

\lstset{style=codestyle}

\titulo{Avaliação 3 -- Projeto RTL do Algoritmo de Máximo Divisor Comum}
\autor{Fabricio Roskowski\\Klaus Dieter Kupper}
\orientador{Prof. Dr. Cesar Albenes Zeferino\\Prof. Dr. Douglas Rossi Melo}
\local{Itajaí -- SC}
\data{Novembro de 2025}
\instituicao{Universidade do Vale do Itajaí\\Programa de Pós-Graduação em Computação Aplicada\\Disciplina: Projeto de Sistemas Digitais}
\preambulo{Relatório técnico apresentado como requisito parcial para a Avaliação 3 da disciplina Projeto de Sistemas Digitais, abordando o desenvolvimento RTL de um sistema digital para cálculo do Máximo Divisor Comum (MDC).}

\makeindex

\begin{document}
\selectlanguage{brazil}
\frenchspacing
\OnehalfSpacing

\imprimircapa
\imprimirfolhaderosto*

\begin{resumo}
Este relatório descreve o desenvolvimento de um sistema digital baseado em VHDL para cálculo do Máximo Divisor Comum (MDC) entre dois operandos de 8 bits, seguindo o método de projeto RTL discutido em sala de aula. Partindo do algoritmo em linguagem C, foi derivada uma máquina de estados de alto nível, identificadas as partições de controle e caminho de dados, definidos os sinais internos com a convenção de nomenclatura solicitada e implementados os blocos VHDL correspondentes (topo, controle, caminho de dados e testbench). São apresentados os diagramas conceituais, a descrição do fluxo de dados, o planejamento de verificação com ModelSim e as orientações para síntese no Quartus Prime. Resultados de simulação e métricas de síntese serão incorporados após a execução dos experimentos.

\vspace{\onelineskip}
\noindent\textbf{Palavras-chave}: VHDL, RTL, máquina de estados, MDC.
\end{resumo}

\tableofcontents*
\textual

\chapter{Introdução}
A Avaliação 3 da disciplina Projeto de Sistemas Digitais tem como objetivo consolidar o método de projeto RTL por meio da implementação de um algoritmo não trivial na linguagem VHDL. O algoritmo selecionado calcula o Máximo Divisor Comum (MDC) de dois operandos inteiros sem sinal de 8 bits utilizando o método de subtrações sucessivas. A especificação exige a decomposição em controle (máquina de estados) e caminho de dados, a adoção de uma convenção de nomenclatura padronizada, a verificação funcional por meio de testbenches no ModelSim e a síntese do sistema no Quartus Prime para obtenção de métricas de custo e desempenho.

Este documento registra o processo completo de desenvolvimento: interpretação dos requisitos, construção da máquina de estados de alto nível a partir do algoritmo em C, identificação dos blocos de controle e processamento, definição dos sinais internos, implementação em VHDL no estilo recomendado pela Xilinx/AMD e preparação do ambiente de verificação. As seções seguintes também deixam espaço reservado para a inclusão de evidências experimentais (formas de onda e diagramas RTL) e resultados de síntese, a serem coletados posteriormente.

\chapter{Especificação do Sistema}

\section{Resumo do Enunciado}
O sistema deve receber duas entradas de dados (\code{i_x} e \code{i_y}), cada uma com 8 bits, e produzir em \code{o_d} o MDC entre esses operandos. O fluxo é controlado por uma entrada de comando \code{i_go} e uma saída de status \code{o_rdy}. Enquanto \code{o_rdy = 1}, o circuito está pronto para receber um novo par de operandos; quando \code{i_go} é ativado, o circuito captura os dados, executa o algoritmo e devolve \code{o_rdy = 1} apenas após o resultado ser disponibilizado em \code{o_d}. A implementação deve ser hierárquica, com módulos separados para topo, controle e caminho de dados, seguindo a convenção de nomes indicada no enunciado (Tabela~\ref{tab:interface}).

\section{Interface e Convenção de Sinais}
A Tabela~\ref{tab:interface} apresenta um resumo das interfaces externas, enquanto a Tabela~\ref{tab:sinais-internos} registra os principais sinais internos que emergiram do processo de particionamento.

\begin{table}[H]
  \centering
\caption{Interface externa do sistema \code{mdc_top}.}
  \label{tab:interface}
  \begin{tabular}{>{\bfseries}p{2.5cm}p{3cm}p{6cm}}
    \toprule
    Identificador & Tipo & Descrição \\
    \midrule
    \code{i_CLK} & Entrada (1 bit) & Clock principal do sistema. \\
    \code{i_RSTn} & Entrada (1 bit) & Reset assíncrono ativo em nível baixo. \\
    \code{i_GO} & Entrada (1 bit) & Pulso que inicia o processamento. \\
    \code{i_X} & Entrada (8 bits) & Operando A sem sinal. \\
    \code{i_Y} & Entrada (8 bits) & Operando B sem sinal. \\
    \code{o_D} & Saída (8 bits) & Resultado do MDC após a convergência. \\
    \code{o_RDY} & Saída (1 bit) & Indica se o módulo está pronto para nova transação. \\
    \bottomrule
  \end{tabular}
\end{table}

\begin{table}[H]
  \centering
  \caption{Sinais internos definidos no particionamento controle/caminho de dados.}
  \label{tab:sinais-internos}
  \begin{tabular}{>{\bfseries}p{3cm}p{3cm}p{6cm}}
    \toprule
    Sinal & Origem/Destino & Função \\
    \midrule
    \code{w_LD_X} & Controle $\rightarrow$ Datapath & Habilita a carga do registrador \code{r_X}. \\
    \code{w_LD_Y} & Controle $\rightarrow$ Datapath & Habilita a carga do registrador \code{r_Y}. \\
    \code{w_SEL_X} & Controle $\rightarrow$ Datapath & Seleciona entre entrada externa e resultado de subtração para \code{r_X}. \\
    \code{w_SEL_Y} & Controle $\rightarrow$ Datapath & Seleciona entre entrada externa e resultado de subtração para \code{r_Y}. \\
    \code{w_X_EQ_Y} & Datapath $\rightarrow$ Controle & Valor lógico indicando igualdade entre \code{r_X} e \code{r_Y}. \\
    \code{w_X_LT_Y} & Datapath $\rightarrow$ Controle & Valor lógico indicando se \code{r_X < r_Y}. \\
    \code{r_X}, \code{r_Y} & Datapath & Registradores que armazenam os valores intermediários do algoritmo. \\
    \bottomrule
  \end{tabular}
\end{table}

\chapter{Metodologia de Projeto RTL}

\section{Algoritmo de Referência em C}
O algoritmo de MDC utilizado como ponto de partida segue o método de subtrações sucessivas, incluindo a lógica de handshake com os sinais \code{i_go} e \code{o_rdy}. O código completo é apresentado na Listagem~\ref{lst:c_mdc}.

\begin{lstlisting}[language=C, caption={Função de referência em C para o cálculo do MDC com handshake}, label={lst:c_mdc}]
#include <stdint.h>
#include <stdbool.h>

static bool i_go = false;
static uint8_t i_x = 0U;
static uint8_t i_y = 0U;

static uint8_t gcd_subtractive(uint8_t x, uint8_t y)
{
  while (x != y) {
    if (x < y) {
      y -= x;
    } else {
      x -= y;
    }
  }
  return x;
}

int main(void)
{
  uint8_t o_d = 0U;
  bool o_rdy = true;

  while (true) {
    while (!i_go) {
      /* idle */
    }

    o_rdy = false;
    o_d = gcd_subtractive(i_x, i_y);
    o_rdy = true;
  }
}
\end{lstlisting}

\section{Derivação da Máquina de Estados}
A partir do algoritmo, identificaram-se seis estados para o controle: \code{s_IDLE}, \code{s_LOAD}, \code{s_COMPARE}, \code{s_SUB_X}, \code{s_SUB_Y} e \code{s_DONE}. A Figura~\ref{fig:fsm-states} registra o diagrama de estados desenhado manualmente (imagem fornecida em \code{media/states.jpg}), enquanto a Tabela~\ref{tab:fsm-transicoes} detalha as transições e saídas relevantes.

\begin{figure}[H]
  \centering
  \includeplaceholder{media/states.jpg}{Reservar espaço para diagrama de estados da FSM}
  \caption{Diagrama preliminar da máquina de estados de controle.}
  \label{fig:fsm-states}
\end{figure}

\begin{table}[H]
  \centering
  \caption{Tabela de transições da FSM de controle.}
  \label{tab:fsm-transicoes}
  \begin{tabular}{>{\bfseries}p{2.5cm}p{3.5cm}p{2.5cm}p{2.5cm}}
    \toprule
    Estado & Ação principal & Condição & Próximo estado \\
    \midrule
    \code{s_IDLE} & \code{o_RDY = 1} & \code{i_GO = 1} & \code{s_LOAD} \\
                  &                  & caso contrário & \code{s_IDLE} \\
    \code{s_LOAD} & Carrega \code{r_X} e \code{r_Y} & -- & \code{s_COMPARE} \\
    \code{s_COMPARE} & Avalia \code{r_X} e \code{r_Y} & \code{w_X_EQ_Y = 1} & \code{s_DONE} \\
                     &                                & \code{w_X_LT_Y = 1} & \code{s_SUB_Y} \\
                     &                                & demais casos & \code{s_SUB_X} \\
    \code{s_SUB_X} & Atualiza \code{r_X} com \code{r_X - r_Y} & -- & \code{s_COMPARE} \\
    \code{s_SUB_Y} & Atualiza \code{r_Y} com \code{r_Y - r_X} & -- & \code{s_COMPARE} \\
    \code{s_DONE} & \code{o_RDY = 1}, mantém \code{o_D} & \code{i_GO = 0} & \code{s_IDLE} \\
    \bottomrule
  \end{tabular}
\end{table}

\section{Particionamento do Caminho de Dados}
O caminho de dados contém dois registradores (\code{r_X} e \code{r_Y}), dois blocos de subtração e multiplexadores de entrada. Os subprodutos das subtrações (\code{w_SUB_XY = r_X - r_Y} e \code{w_SUB_YX = r_Y - r_X}) são selecionados pelo controle para atualizar o registrador apropriado, evitando underflow por meio da lógica de decisão da FSM. As comparações de igualdade (\code{w_X_EQ_Y}) e menor (\code{w_X_LT_Y}) são calculadas continuamente e alimentam o bloco de controle.

\chapter{Arquitetura RTL Proposta}

\section{Visão Geral}
A Figura~\ref{fig:rtl-diagram} apresenta o diagrama RTL em alto nível, destacando a interação entre o bloco de controle (\code{mdc_control}) e o caminho de dados (\code{mdc_datapath}). As setas são etiquetadas com os sinais definidos na Tabela~\ref{tab:sinais-internos}.

\begin{figure}[H]
  \centering
  \includeplaceholder{media/mdc_diagram_example.png}{Reservar espaço para o diagrama RTL (controle + caminho de dados)}
  \caption{Particionamento em controle e caminho de dados do sistema MDC.}
  \label{fig:rtl-diagram}
\end{figure}

\section{Bloco de Controle}
O bloco \code{mdc_control} é responsável por sequenciar o algoritmo, garantindo que apenas operações válidas de subtração sejam realizadas. O uso de estados explícitos facilita a inclusão de novas condições de parada (por exemplo, detecção de operandos nulos) sem alterar o caminho de dados. As saídas \code{w_LD_X}, \code{w_LD_Y}, \code{w_SEL_X} e \code{w_SEL_Y} são exclusivamente geradas pelo controle e utilizadas como sinais de habilitação e seleção dentro do datapath.

\section{Caminho de Dados}
O caminho de dados mantém os operandos em registradores e atualiza seus valores com base na decisão do controle. As subtrações são realizadas em paralelo, permitindo que o controlador selecione o resultado correto com atraso mínimo. O resultado final é exposto diretamente em \code{o_result}, que é repassado ao módulo topo.

\chapter{Implementação em VHDL}
Os códigos foram escritos seguindo o guia de estilo da avaliação: palavras reservadas em minúsculas, sinais em maiúsculas com prefixos (\code{i_}, \code{o_}, \code{w_}, \code{r_}), tabulação de dois espaços e comentários em inglês para evitar problemas de codificação.

\section{Módulo Topo}
A Listagem~\ref{lst:vhdl_top} mostra o módulo topo \code{mdc_top}, responsável por conectar os blocos de controle e caminho de dados e disponibilizar a interface externa.

\begin{lstlisting}[language=VHDL, caption={Módulo topo do sistema MDC}, label={lst:vhdl_top}]
library ieee;
use ieee.std_logic_1164.all;

entity mdc_top is
  port (
    i_clk  : in  std_logic;
    i_rstn : in  std_logic;
    i_go   : in  std_logic;
    i_x    : in  std_logic_vector(7 downto 0);
    i_y    : in  std_logic_vector(7 downto 0);
    o_d    : out std_logic_vector(7 downto 0);
    o_rdy  : out std_logic
  );
end entity mdc_top;

architecture rtl of mdc_top is
  signal w_LD_X   : std_logic;
  signal w_LD_Y   : std_logic;
  signal w_SEL_X  : std_logic;
  signal w_SEL_Y  : std_logic;
  signal w_X_EQ_Y : std_logic;
  signal w_X_LT_Y : std_logic;
  signal w_RESULT : std_logic_vector(7 downto 0);
begin

  u_datapath : entity work.mdc_datapath
    port map (
      i_clk    => i_clk,
      i_rstn   => i_rstn,
      i_x      => i_x,
      i_y      => i_y,
      i_ld_x   => w_LD_X,
      i_ld_y   => w_LD_Y,
      i_sel_x  => w_SEL_X,
      i_sel_y  => w_SEL_Y,
      o_x_eq_y => w_X_EQ_Y,
      o_x_lt_y => w_X_LT_Y,
      o_result => w_RESULT
    );

  u_control : entity work.mdc_control
    port map (
      i_clk    => i_clk,
      i_rstn   => i_rstn,
      i_go     => i_go,
      i_x_eq_y => w_X_EQ_Y,
      i_x_lt_y => w_X_LT_Y,
      o_ld_x   => w_LD_X,
      o_ld_y   => w_LD_Y,
      o_sel_x  => w_SEL_X,
      o_sel_y  => w_SEL_Y,
      o_rdy    => o_rdy
    );

  o_d <= w_RESULT;

end architecture rtl;
\end{lstlisting}

\section{Máquina de Controle}
A Listagem~\ref{lst:vhdl_control} apresenta a implementação da FSM derivada na Tabela~\ref{tab:fsm-transicoes}. Os estados são codificados em um tipo enumerado \code{t_STATE}, e os sinais de controle seguem diretamente a tabela de transições.

\begin{lstlisting}[language=VHDL, caption={Implementação da FSM de controle}, label={lst:vhdl_control}]
library ieee;
use ieee.std_logic_1164.all;

entity mdc_control is
  port (
    i_clk    : in  std_logic;
    i_rstn   : in  std_logic;
    i_go     : in  std_logic;
    i_x_eq_y : in  std_logic;
    i_x_lt_y : in  std_logic;
    o_ld_x   : out std_logic;
    o_ld_y   : out std_logic;
    o_sel_x  : out std_logic;
    o_sel_y  : out std_logic;
    o_rdy    : out std_logic
  );
end entity mdc_control;

architecture rtl of mdc_control is
  type t_STATE is (s_IDLE, s_LOAD, s_COMPARE, s_SUB_X, s_SUB_Y, s_DONE);
  signal r_STATE      : t_STATE;
  signal w_NEXT_STATE : t_STATE;

  signal w_LD_X   : std_logic;
  signal w_LD_Y   : std_logic;
  signal w_SEL_X  : std_logic;
  signal w_SEL_Y  : std_logic;
  signal w_RDY    : std_logic;
begin

  p_STATE_REG : process(i_clk, i_rstn)
  begin
    if i_rstn = '0' then
      r_STATE <= s_IDLE;
    elsif rising_edge(i_clk) then
      r_STATE <= w_NEXT_STATE;
    end if;
  end process p_STATE_REG;

  p_STATE_NEXT : process(r_STATE, i_go, i_x_eq_y, i_x_lt_y)
  begin
    w_NEXT_STATE <= r_STATE;
    case r_STATE is
      when s_IDLE =>
        if i_go = '1' then
          w_NEXT_STATE <= s_LOAD;
        end if;
      when s_LOAD =>
        w_NEXT_STATE <= s_COMPARE;
      when s_COMPARE =>
        if i_x_eq_y = '1' then
          w_NEXT_STATE <= s_DONE;
        elsif i_x_lt_y = '1' then
          w_NEXT_STATE <= s_SUB_Y;
        else
          w_NEXT_STATE <= s_SUB_X;
        end if;
      when s_SUB_X =>
        w_NEXT_STATE <= s_COMPARE;
      when s_SUB_Y =>
        w_NEXT_STATE <= s_COMPARE;
      when s_DONE =>
        if i_go = '0' then
          w_NEXT_STATE <= s_IDLE;
        end if;
    end case;
  end process p_STATE_NEXT;

  p_OUTPUTS : process(r_STATE)
  begin
    w_LD_X  <= '0';
    w_LD_Y  <= '0';
    w_SEL_X <= '0';
    w_SEL_Y <= '0';
    w_RDY   <= '0';

    case r_STATE is
      when s_IDLE =>
        w_RDY <= '1';
      when s_LOAD =>
        w_LD_X  <= '1';
        w_LD_Y  <= '1';
      when s_COMPARE =>
        null;
      when s_SUB_X =>
        w_LD_X  <= '1';
        w_SEL_X <= '1';
      when s_SUB_Y =>
        w_LD_Y  <= '1';
        w_SEL_Y <= '1';
      when s_DONE =>
        w_RDY <= '1';
    end case;
  end process p_OUTPUTS;

  o_ld_x  <= w_LD_X;
  o_ld_y  <= w_LD_Y;
  o_sel_x <= w_SEL_X;
  o_sel_y <= w_SEL_Y;
  o_rdy   <= w_RDY;

end architecture rtl;
\end{lstlisting}

\section{Caminho de Dados}
A Listagem~\ref{lst:vhdl_datapath} contém a descrição do caminho de dados. Os cálculos são realizados com a biblioteca \code{numeric_std}, garantindo aritmética sem sinal consistente com a especificação.

\begin{lstlisting}[language=VHDL, caption={Descrição do caminho de dados}, label={lst:vhdl_datapath}]
library ieee;
use ieee.std_logic_1164.all;
use ieee.numeric_std.all;

entity mdc_datapath is
  port (
    i_clk    : in  std_logic;
    i_rstn   : in  std_logic;
    i_x      : in  std_logic_vector(7 downto 0);
    i_y      : in  std_logic_vector(7 downto 0);
    i_ld_x   : in  std_logic;
    i_ld_y   : in  std_logic;
    i_sel_x  : in  std_logic;
    i_sel_y  : in  std_logic;
    o_x_eq_y : out std_logic;
    o_x_lt_y : out std_logic;
    o_result : out std_logic_vector(7 downto 0)
  );
end entity mdc_datapath;

architecture rtl of mdc_datapath is
  signal r_X      : std_logic_vector(7 downto 0);
  signal r_Y      : std_logic_vector(7 downto 0);
  signal w_NEXT_X : std_logic_vector(7 downto 0);
  signal w_NEXT_Y : std_logic_vector(7 downto 0);
  signal w_SUB_XY : std_logic_vector(7 downto 0);
  signal w_SUB_YX : std_logic_vector(7 downto 0);
begin

  w_SUB_XY <= std_logic_vector(unsigned(r_X) - unsigned(r_Y));
  w_SUB_YX <= std_logic_vector(unsigned(r_Y) - unsigned(r_X));

  with i_sel_x select
    w_NEXT_X <= i_x      when '0',
                w_SUB_XY when others;

  with i_sel_y select
    w_NEXT_Y <= i_y      when '0',
                w_SUB_YX when others;

  p_REGISTERS : process(i_clk, i_rstn)
  begin
    if i_rstn = '0' then
      r_X <= (others => '0');
      r_Y <= (others => '0');
    elsif rising_edge(i_clk) then
      if i_ld_x = '1' then
        r_X <= w_NEXT_X;
      end if;
      if i_ld_y = '1' then
        r_Y <= w_NEXT_Y;
      end if;
    end if;
  end process p_REGISTERS;

  o_x_eq_y <= '1' when r_X = r_Y else '0';
  o_x_lt_y <= '1' when unsigned(r_X) < unsigned(r_Y) else '0';
  o_result <= r_X;

end architecture rtl;
\end{lstlisting}

\section{Testbench}
O testbench da Listagem~\ref{lst:vhdl_tb} instancia o módulo topo, gera o clock, aplica vetores de teste e aguarda o término de cada transação antes de iniciar a próxima. As assertivas de verificação podem ser ampliadas conforme necessário.

\begin{lstlisting}[language=VHDL, caption={Testbench funcional para o sistema MDC}, label={lst:vhdl_tb}]
library ieee;
use ieee.std_logic_1164.all;
use ieee.numeric_std.all;

entity mdc_tb is
end entity mdc_tb;

architecture sim of mdc_tb is
  constant c_CLK_PERIOD : time := 10 ns;

  signal w_CLK  : std_logic := '0';
  signal w_RSTn : std_logic := '0';
  signal w_GO   : std_logic := '0';
  signal w_X    : std_logic_vector(7 downto 0) := (others => '0');
  signal w_Y    : std_logic_vector(7 downto 0) := (others => '0');
  signal w_D    : std_logic_vector(7 downto 0);
  signal w_RDY  : std_logic;
begin

  u_dut : entity work.mdc_top
    port map (
      i_clk  => w_CLK,
      i_rstn => w_RSTn,
      i_go   => w_GO,
      i_x    => w_X,
      i_y    => w_Y,
      o_d    => w_D,
      o_rdy  => w_RDY
    );

  p_CLK : process
  begin
    w_CLK <= '0';
    wait for c_CLK_PERIOD / 2;
    w_CLK <= '1';
    wait for c_CLK_PERIOD / 2;
  end process p_CLK;

  p_STIMULUS : process
  begin
    w_RSTn <= '0';
    wait for 3 * c_CLK_PERIOD;
    w_RSTn <= '1';
    wait for c_CLK_PERIOD;

    w_X <= std_logic_vector(to_unsigned(108, 8));
    w_Y <= std_logic_vector(to_unsigned(84, 8));
    wait for c_CLK_PERIOD;
    w_GO <= '1';
    wait for c_CLK_PERIOD;
    w_GO <= '0';

    wait until w_RDY = '1';
    wait for 2 * c_CLK_PERIOD;

    w_X <= std_logic_vector(to_unsigned(27, 8));
    w_Y <= std_logic_vector(to_unsigned(18, 8));
    wait for c_CLK_PERIOD;
    w_GO <= '1';
    wait for c_CLK_PERIOD;
    w_GO <= '0';

    wait for 60 * c_CLK_PERIOD;
    wait;
  end process p_STIMULUS;

end architecture sim;
\end{lstlisting}

\chapter{Plano de Validação e Síntese}

\section{Validação Funcional}
A simulação será conduzida no ModelSim, utilizando o testbench descrito na Seção anterior. Estão planejadas as seguintes etapas:
\begin{enumerate}[label=\alph*)]
  \item Compilação dos arquivos VHDL (\code{mdc_datapath.vhd}, \code{mdc_control.vhd}, \code{mdc_top.vhd}, \code{mdc_tb.vhd}) com suporte a VHDL-2008.
  \item Execução do testbench até que \code{o_RDY} retorne ao nível lógico alto, capturando as formas de onda relevantes (\code{i_go}, \code{i_x}, \code{i_y}, \code{o_d}, \code{o_rdy}, \code{w_X_EQ_Y}, \code{w_X_LT_Y}, \code{r_X}, \code{r_Y}).
  \item Inclusão das figuras de forma de onda no relatório (Figura~\ref{fig:simulation-waveforms}) com marcação dos pontos de interesse.
\end{enumerate}

\begin{figure}[H]
  \centering
  \includeplaceholder{media/mdc_waveforms_example.png}{Reservar espaço para as formas de onda de simulação}
  \caption{Espaço reservado para as formas de onda obtidas no ModelSim.}
  \label{fig:simulation-waveforms}
\end{figure}

\section{Síntese e Análise RTL}
A síntese será realizada no Quartus Prime, considerando a mesma hierarquia adotada na simulação. Após a compilação, será exportado o diagrama RTL para inclusão na Figura~\ref{fig:quartus-rtl}, bem como coletadas as métricas de utilização de recursos e frequência máxima (Tabela~\ref{tab:resource-usage}).

\begin{figure}[H]
  \centering
  \includeplaceholder{media/mdc_quartus_rtl_example.png}{Reservar espaço para o diagrama RTL gerado pelo Quartus}
  \caption{Espaço reservado para o diagrama RTL gerado pelo Quartus Prime.}
  \label{fig:quartus-rtl}
\end{figure}

\begin{table}[H]
  \centering
  \caption{Template para registrar os custos e o desempenho após a síntese.}
  \label{tab:resource-usage}
  \begin{tabular}{>{\bfseries}p{6cm}p{3cm}p{3cm}}
    \toprule
    Métrica & Valor & Observações \\
    \midrule
    Número de LCs & & \\
    Número de LUTs & & \\
    Número de FFs & & \\
    Frequência máxima (Fmax) & & \\
    \bottomrule
  \end{tabular}
\end{table}

\chapter{Conclusões Parciais}
O projeto encontra-se em estágio funcional, com os blocos de controle e datapath implementados e integrados. A metodologia de projeto RTL foi seguida rigorosamente: partiu-se do algoritmo em C, derivou-se uma FSM de alto nível, identificaram-se os sinais necessários e codificaram-se os módulos em VHDL com a convenção estabelecida. As próximas atividades concentram-se na coleta de evidências experimentais (simulações e síntese) e no refinamento do relatório com base nesses resultados.

\postextual
\phantomsection
\addcontentsline{toc}{chapter}{Referências}
\begin{thebibliography}{9}
\bibitem{vhdl-guide} Xilinx. \textit{VHDL Coding Style Guidelines}. AMD Adaptive Computing Documentation, 2023.
\bibitem{quartus} Intel. \textit{Quartus Prime Standard Edition User Guide}. 2024.
\bibitem{modelsim} Siemens EDA. \textit{ModelSim Simulation User Manual}. 2024.
\end{thebibliography}

\end{document}


